% GENERATED FILE. Edit canonical lessons, then run npm run generate:books.
% canonical-source-hash: fnv1a64:ccdd16872b0c24b2
% canonical-lessons: JA-C152-shitaku, JA-C152-isogu, JA-C152-dekakeru, JA-C152-wasureru, JA-C152-okureru

\chapter{Hurrying to Get Out of the House}
\label{ch:ja-hurrying-to-get-out-of-the-house}
\hlchaptermodality{152}

\begin{chapteropening}
\textbf{By the end of this chapter:} \emph{I can say how I get ready and leave the house.}

Complete the last lesson of chapter 152: I can say how I get ready and leave the house.
\end{chapteropening}

\section[\texorpdfstring{\ja{したく} (shitaku)}{したく (shitaku)}]{\ja{したく} (shitaku) — preparations; getting things ready}

\label{lesson:JA-C152-shitaku}



\begin{quote}
Before the new one: say the Japanese for a sleeve, then the Japanese for to suit.
\end{quote}

\subsection*{You'll want to know: \ja{したく}}
\textbf{\ja{したく}} — \emph{shitaku} — \textquotedblleft{}preparations; getting things ready\textquotedblright{}.

\begin{quote}
\textbf{\ja{あさの} \ja{したくは} \ja{かんたんです。}} — \emph{asa no shitaku wa kantan desu} — \textquotedblleft{}Getting ready in the morning is simple.\textquotedblright{}
\end{quote}

\begin{cousinweb}
\textbf{\ja{したく}} (\emph{shitaku}) is Sino-Japanese (kango): two Chinese characters read with their on'yomi, readings borrowed from Chinese. They are \emph{shi}, \textquotedblleft{}branch, support\textquotedblright{} (Middle Chinese \emph{tsye}) and \emph{taku}, \textquotedblleft{}degrees, consider\textquotedblright{} (\emph{dak}) (KANJIDIC2; Middle Chinese in Baxter's transcription).
\end{cousinweb}

\begin{grammarlens}[title={What you've built}]
One more word for everyday life at A2.
\end{grammarlens}

\subsection*{Your turn}
\begin{itemize}
\raggedright
  \item \emph{Say it:} \emph{shitaku}
  \item \emph{Say it:} \emph{shitaku}, once more
  \item \emph{Say it:} \emph{asa no shitaku wa kantan desu} — Getting ready in the morning is simple.
  \item \emph{Recall:} say the Japanese for a sleeve, then the Japanese for to suit, then say \emph{shitaku} again
\end{itemize}

\begin{tcolorbox}[breakable,colback=teal!4,colframe=teal!35!black,title={Before you move on}]
What does \ja{したく} mean? (\textquotedblleft{}Preparations; getting things ready\textquotedblright{}.) Say it once more.
\end{tcolorbox}
\section[\texorpdfstring{\ja{いそぐ} (isogu)}{いそぐ (isogu)}]{\ja{いそぐ} (isogu) — to hurry}

\label{lesson:JA-C152-isogu}



\begin{quote}
Before the new one: say the Japanese for to suit, then the Japanese for preparations.
\end{quote}

\subsection*{You'll want to know: \ja{いそぐ}}
\textbf{\ja{いそぐ}} — \emph{isogu} — \textquotedblleft{}to hurry\textquotedblright{}.

\begin{quote}
\textbf{\ja{あさは} \ja{いつも} \ja{いそぎます。}} — \emph{asa wa itsumo isogimasu} — \textquotedblleft{}In the morning I am always in a hurry.\textquotedblright{}
\end{quote}

Polite \textbf{\ja{いそぎます}}, \emph{isogimasu}; a godan verb.

\begin{cousinweb}
\textbf{\ja{いそぐ}} (\emph{isogu}) is a native Japanese word (wago), inherited rather than borrowed from Chinese; its kanji was attached later as a kun'yomi, and its older origin is not settled (JMdict, KANJIDIC2).
\end{cousinweb}

\begin{grammarlens}[title={What you've built}]
One more word for everyday life at A2.
\end{grammarlens}

\subsection*{Your turn}
\begin{itemize}
\raggedright
  \item \emph{Say it:} \emph{isogu}
  \item \emph{Say it:} \emph{isogu}, once more
  \item \emph{Say it:} \emph{asa wa itsumo isogimasu} — In the morning I am always in a hurry.
  \item \emph{Recall:} say the Japanese for to suit, then the Japanese for preparations, then say \emph{isogu} again
\end{itemize}

\begin{tcolorbox}[breakable,colback=teal!4,colframe=teal!35!black,title={Before you move on}]
What does \ja{いそぐ} mean? (\textquotedblleft{}To hurry\textquotedblright{}.) Say it once more.
\end{tcolorbox}
\section[\texorpdfstring{\ja{でかける} (dekakeru)}{でかける (dekakeru)}]{\ja{でかける} (dekakeru) — to go out; to set off}

\label{lesson:JA-C152-dekakeru}



\begin{quote}
Before the new one: say the Japanese for preparations, then the Japanese for to hurry.
\end{quote}

\subsection*{You'll want to know: \ja{でかける}}
\textbf{\ja{でかける}} — \emph{dekakeru} — \textquotedblleft{}to go out; to set off\textquotedblright{}.

\begin{quote}
\textbf{\ja{ちちは} \ja{まいあさ} \ja{でかけます。}} — \emph{chichi wa maiasa dekakemasu} — \textquotedblleft{}My father goes out every morning.\textquotedblright{}
\end{quote}

Polite \textbf{\ja{でかけます}}, \emph{dekakemasu}; an ichidan verb.

\begin{cousinweb}
\textbf{\ja{でかける}} (\emph{dekakeru}) is native Japanese, built from \emph{de}, \textquotedblleft{}going out\textquotedblright{}, the stem of \emph{deru}, and \emph{kakeru}, \textquotedblleft{}to set about\textquotedblright{}, both inherited rather than borrowed (JMdict, KANJIDIC2).
\end{cousinweb}

\begin{grammarlens}[title={What you've built}]
One more word for everyday life at A2.
\end{grammarlens}

\subsection*{Your turn}
\begin{itemize}
\raggedright
  \item \emph{Say it:} \emph{dekakeru}
  \item \emph{Say it:} \emph{dekakeru}, once more
  \item \emph{Say it:} \emph{chichi wa maiasa dekakemasu} — My father goes out every morning.
  \item \emph{Recall:} say the Japanese for preparations, then the Japanese for to hurry, then say \emph{dekakeru} again
\end{itemize}

\begin{tcolorbox}[breakable,colback=teal!4,colframe=teal!35!black,title={Before you move on}]
What does \ja{でかける} mean? (\textquotedblleft{}To go out; to set off\textquotedblright{}.) Say it once more.
\end{tcolorbox}
\section[\texorpdfstring{\ja{わすれる} (wasureru)}{わすれる (wasureru)}]{\ja{わすれる} (wasureru) — to forget; to leave behind}

\label{lesson:JA-C152-wasureru}



\begin{quote}
Before the new one: say the Japanese for to hurry, then the Japanese for to go out.
\end{quote}

\subsection*{You'll want to know: \ja{わすれる}}
\textbf{\ja{わすれる}} — \emph{wasureru} — \textquotedblleft{}to forget; to leave behind\textquotedblright{}.

\begin{quote}
\textbf{\ja{ちちは} \ja{よく} \ja{かさを} \ja{わすれます。}} — \emph{chichi wa yoku kasa o wasuremasu} — \textquotedblleft{}My father often forgets his umbrella.\textquotedblright{}
\end{quote}

Polite \textbf{\ja{わすれます}}, \emph{wasuremasu}; an ichidan verb.

\begin{cousinweb}
\textbf{\ja{わすれる}} (\emph{wasureru}) is a native Japanese word (wago), inherited rather than borrowed from Chinese; its kanji was attached later as a kun'yomi, and its older origin is not settled (JMdict, KANJIDIC2).
\end{cousinweb}

\begin{grammarlens}[title={What you've built}]
One more word for everyday life at A2.
\end{grammarlens}

\subsection*{Your turn}
\begin{itemize}
\raggedright
  \item \emph{Say it:} \emph{wasureru}
  \item \emph{Say it:} \emph{wasureru}, once more
  \item \emph{Say it:} \emph{chichi wa yoku kasa o wasuremasu} — My father often forgets his umbrella.
  \item \emph{Recall:} say the Japanese for to hurry, then the Japanese for to go out, then say \emph{wasureru} again
\end{itemize}

\begin{tcolorbox}[breakable,colback=teal!4,colframe=teal!35!black,title={Before you move on}]
What does \ja{わすれる} mean? (\textquotedblleft{}To forget; to leave behind\textquotedblright{}.) Say it once more.
\end{tcolorbox}
\section[\texorpdfstring{\ja{おくれる} (okureru)}{おくれる (okureru)}]{\ja{おくれる} (okureru) — to be late}

\label{lesson:JA-C152-okureru}



\begin{quote}
Before the new one: say the Japanese for to go out, then the Japanese for to forget.
\end{quote}

\subsection*{You'll want to know: \ja{おくれる}}
\textbf{\ja{おくれる}} — \emph{okureru} — \textquotedblleft{}to be late\textquotedblright{}.

\begin{quote}
\textbf{\ja{いもうとは} \ja{よく} \ja{がっこうに} \ja{おくれます。}} — \emph{imōto wa yoku gakkō ni okuremasu} — \textquotedblleft{}My younger sister is often late for school.\textquotedblright{}
\end{quote}

Polite \textbf{\ja{おくれます}}, \emph{okuremasu}; an ichidan verb.

\begin{cousinweb}
\textbf{\ja{おくれる}} (\emph{okureru}) is a native Japanese word (wago), inherited rather than borrowed from Chinese; its kanji was attached later as a kun'yomi, and its older origin is not settled (JMdict, KANJIDIC2).
\end{cousinweb}

\begin{grammarlens}[title={What you've built}]
One more word for everyday life at A2.
\end{grammarlens}

\subsection*{Your turn}
\begin{itemize}
\raggedright
  \item \emph{Say it:} \emph{okureru}
  \item \emph{Say it:} \emph{okureru}, once more
  \item \emph{Say it:} \emph{imōto wa yoku gakkō ni okuremasu} — My younger sister is often late for school.
  \item \emph{Recall:} say the Japanese for to go out, then the Japanese for to forget, then say \emph{okureru} again
\end{itemize}

\begin{tcolorbox}[breakable,colback=teal!4,colframe=teal!35!black,title={Before you move on}]
What does \ja{おくれる} mean? (\textquotedblleft{}To be late\textquotedblright{}.) Say it once more.
\end{tcolorbox}
